% Part 2 chapter carrying "Electroweak Symmetry Breaking and the Matter Sector" (Mar 2026, revised Oct 2026), read in full 2026-10-02 (PDF text
% 207 lines; ledger docs/book/read_ledgers/LEDGER_G8_35_Electroweak.md). Corrections applied: PAPER_ERRATA EW1, EW2, SP5/EN3 (DESI combination dropped).
% Numbers: docs/verification/scripts/verify_entanglement_electroweak.py (W1-W12). Check: docs/verification/particle/ELECTROWEAK_CHECK.md.
% Already carried and only referenced here: the matter-sector criterion, E(a) at the electroweak epoch, m_gamma bound and beta_m composition
% (Chapter ch:higgsrecord, Section sec:ew:betam); eta and the Sakharov conditions (ch:baryon); before/after the transition and the arrow of time (ch:time, Part 5).
\chapter{Electroweak symmetry breaking and the matter sector}\label{ch:electroweak}

\section*{The place}
The result of this chapter is that a matter sector able to decohere gravitationally exists from electroweak symmetry breaking, and the
photon stays outside it by an exact gauge symmetry. This chapter gives the particle physics behind that statement: which interactions bind and how,
which symmetries the weak interaction breaks, what the transition is and when it happens, and what it does not establish. Part~5 treats the same
epoch from the side of time (Chapter~\ref{ch:time}).

\section{Which interactions bind}\label{sec:ew:bind}
IAM's Law comes first: every irreversible record pays $\kB T\ln2$ per bit to the nearest encoding surface. For bound matter the virial partition then
sets how much is written (Chapters~\ref{ch:virial_law} and~\ref{ch:virial}). The partition depends on the form of the potential. For a potential
homogeneous of degree $k$, $V(\lambda r)=\lambda^kV(r)$, the time-averaged virial theorem gives~\cite{Clausius1870}
\begin{equation}
2\langle K\rangle=k\,\langle V\rangle.\label{eq:ew:virial}
\end{equation}
\derived{} For the $1/r$ potentials of gravity and electrostatics, $k=-1$ and $2\langle K\rangle+\langle V\rangle=0$. This holds for atoms and
molecules, for stars and for galaxy clusters, and it is the basis of the coupling $\beta_m=\Omega_m/2$ (Chapter~\ref{ch:virial}).

The strong interaction is different. The static quark potential is well described by the Cornell form~\cite{Eichten1978}
\begin{equation}
V(r)=-\frac43\frac{\alpha_s}{r}+\sigma r,\qquad\sigma\approx0.18\ {\rm GeV}^2=0.91\ {\rm GeV\,fm}^{-1}.
\end{equation}
\calc{} It is Coulomb-like ($k=-1$) at short distance and linear ($k=+1$) at the confinement scale, where Eq.~\eqref{eq:ew:virial} gives
$2\langle K\rangle=+\langle V_{\rm lin}\rangle$. Hadrons are virialised bound states, but of their own potential, not of the $1/r$ form. How much a
hadron writes is not treated. \openprob

The weak interaction forms no bound states. Its range is set by the $W$ mass, $\hbar/(M_Wc)=2.5\times10^{-18}$\,m for $M_W=80.377$\,GeV~\cite{PDG2022}.
\calc{} Its role is to change flavour.

\section{The weak interaction and the discrete symmetries}\label{sec:ew:weak}
The weak interaction is carried by the $W^\pm$ ($80.4$\,GeV) and the $Z^0$ ($91.2$\,GeV)~\cite{PDG2022}. It is the only interaction that changes quark
and lepton flavour. It drives beta decay and the first step of the proton--proton chain, $p+p\to d+e^++\nu_e$. It violates parity
maximally~\cite{Wu1957}. It violates CP through the complex phase of the quark mixing matrix~\cite{Cabibbo1963,KobayashiMaskawa1973}, first seen in
neutral kaon decays~\cite{Christenson1964}. By the CPT theorem it therefore violates time-reversal symmetry in its fundamental laws. \observed

Two points matter for the matter sector.
\begin{enumerate}
\item CP violation is one of the three Sakharov conditions for a baryon asymmetry~\cite{Sakharov1967}. The asymmetry $\eta$ fixes $\Omega_b$, and
with it the baryonic part of $\beta_m$ (Section~\ref{sec:ew:betam} and Chapter~\ref{ch:baryon}). The phase of the quark mixing matrix is too small to produce
the measured $\eta$. In the Standard Model the electroweak transition is a smooth crossover~\cite{Kajantie1996,DOnofrio2016}, not the first-order
transition that electroweak baryogenesis needs. The mechanism that set $\eta$ is unknown. \observed
\item The thermodynamic arrow does not need CP violation. Entropy production, decoherence and Landauer's bound hold for every interaction. What the
electroweak transition adds is not irreversibility but massive particles. \interp
\end{enumerate}

\section{The transition}\label{sec:ew:transition}
For the measured Higgs mass the Standard Model has no electroweak phase transition: lattice studies find a crossover~\cite{Kajantie1996}, centred at
$T_c=159.5\pm1.5$\,GeV~\cite{DOnofrio2016}. \calc{} In the radiation era $t=0.301\,g_*^{-1/2}\,m_P/T^2$ (in $\hbar=c=k_B=1$ units); with
$g_*=106.75$ the crossover falls at $t\approx9\times10^{-12}$\,s. Entropy conservation gives the scale factor there, $a_{\rm EW}=4.9\times10^{-16}$,
so $\ln E(a_{\rm EW})=1-1/a_{\rm EW}=-2.0\times10^{15}$. \calc{} Because the change is smooth, no single temperature marks the moment masses appear;
the statements below refer to the two sides of the crossover.

\paragraph{Above.} The Higgs condensate is negligible. The $W$, $Z$ and the fermions have no rest mass, and their excitations travel on null
geodesics. In the plasma they carry thermal masses. These are a collective effect of the medium that vanishes in vacuum. The criterion
$\Delta\tau>0$ is stated for rest mass; how it applies to quasiparticles in a hot plasma is not treated. \openprob

\paragraph{Below.} The Higgs field takes its vacuum value $v=(\sqrt2G_F)^{-1/2}=246.22$\,GeV. \calc{} The $W^\pm$ and $Z^0$ acquire mass, the charged
fermions acquire $m_f=y_fv/\sqrt2$, and their excitations move on timelike worldlines with $\Delta\tau>0$. By the criterion of
Section~\ref{sec:hr:proper} (a degree of freedom decoheres gravitationally only if it accumulates proper time), the matter sector exists from here on. Dark matter, whose particle nature is unknown, joins it by the same criterion if it
is massive. Neutrinos have mass, shown by their oscillations, so they too move on timelike worldlines; the origin of their mass lies outside the
minimal Standard Model~\cite{PDG2022}. \observed

\paragraph{The photon.} The electromagnetic $U(1)$ survives unbroken, so the photon stays exactly massless and on null geodesics
($m_\gamma<10^{-18}$\,eV~\cite{PDG2024}). Its exemption from writing follows from an established symmetry, not from an added assumption. \derived{}
It does not settle $\Sigma=1$. $\Sigma$ describes how the lensing potential $\Phi+\Psi$ responds to matter. In the framework $\Sigma=1$ follows from
$\delta\varphi=0$ and from the absence of anisotropic stress in the informational sector (Chapter~\ref{ch:theory}). Weak lensing tests it directly.

\section{The sequence of bound structures}\label{sec:ew:sequence}
The transition sits in the standard thermal history as follows. Times are for the Planck 2018 cosmology~\cite{Planck2018VI}; energies are those of the
Standard Model~\cite{PDG2022}.
\begin{center}\small\begin{tabularx}{\linewidth}{>{\raggedright\arraybackslash}p{0.16\linewidth}>{\raggedright\arraybackslash}p{0.12\linewidth}>{\raggedright\arraybackslash}p{0.17\linewidth}>{\raggedright\arraybackslash}X}\toprule
time & temperature & event & consequence\\\midrule
$9\times10^{-12}$\,s & 159.5\,GeV & electroweak crossover & $W$, $Z$, fermions massive; photon massless; matter sector exists\\
$\sim10^{-5}$\,s & $\sim150$\,MeV & QCD confinement & hadrons: bound states of the Cornell potential\\
2--4.5\,min & 0.1--0.07\,MeV & nucleosynthesis & light nuclei; $\Omega_b$ already fixed by $\eta$\\
$3.7\times10^5$\,yr & 0.256\,eV & recombination & atoms: $2\langle K\rangle+\langle V\rangle=0$ (Coulomb)\\
$\sim100$\,Myr ($z\approx30$) & --- & first haloes & gravitational virialisation; $E=9.4\times10^{-14}$\\
13.8\,Gyr & 2.7255\,K & today & $E=1$; $\mu(z{=}0)=1/(1+\beta_m)=0.864$\\\bottomrule
\end{tabularx}\end{center}
\calc{} The matter sector exists from the first row. It writes in quantity only from the fifth, once structure forms.

\section{The composition of $\beta_m$}\label{sec:ew:betam}
The coupling $\beta_m=\Omega_m/2$ (Chapter~\ref{ch:virial}) divides between the two kinds of matter. With the Planck 2018 values $\Omega_m=0.3153$
and $\Omega_b=0.0493$~\cite{Planck2018VI},
\begin{equation}
\beta_m=\frac{\Omega_b}{2}+\frac{\Omega_{\rm dm}}{2}=0.0247+0.1330=0.1577.\label{eq:ew:betam}
\end{equation}
\calc{} The baryonic part, $15.6\,\%$, is set by the baryon asymmetry (Chapter~\ref{ch:baryon}); the dark part, $84.4\,\%$, by the dark-matter
density, which the microwave background fixes at $z\approx1090$ (Chapter~\ref{ch:higgsrecord}). Both are measured inputs; $\beta_m$ adds no
parameter to them. That $\beta_m/\Omega_m=1/2$ for any revision of $\Omega_m$ is the definition, not a test. \derived

\section{The vacuum question}\label{sec:ew:vacuum}
Chapter~\ref{ch:electronmass} leaves open, and Chapter~\ref{ch:time} in Part~\ref{part:5} returns to, whether the settling of the Higgs field into its vacuum is itself a record written at
the temperature of the transition. The particle physics narrows the question. The points of the Higgs vacuum manifold are related by local gauge
transformations. A local gauge symmetry cannot break spontaneously~\cite{Elitzur1975}, so these points are not distinct physical states, and no
gauge-invariant quantity records which one was taken. The gauge-invariant change across the crossover, the growth of $\langle\phi^\dagger\phi\rangle$,
is smooth and selects nothing. A record needs distinct physical outcomes, as in a first-order transition that leaves domains, which the Standard Model
does not provide. The open form of the question is therefore whether any physical outcome is selected at the electroweak epoch. Nothing in the
cosmology depends on it. \openprob

\section{Tests}\label{sec:ew:tests}
\begin{enumerate}
\item $\Sigma=1$ at every redshift: Euclid tomographic weak lensing. A significant $\Sigma\neq1$ falsifies the photon exemption. \prediction
\item $\mu_0=-0.136$, that is $\mu(z{=}0)=0.864$: Euclid clustering. The complete first data release is due in mid-2027. Euclid's published forecasts bracket the test of the IAM $\mu(z)$ from about $0.3\sigma$ (conservative scale cuts) to about $7\sigma$ (smallest
scales modelled), and a forecast made with the IAM $\mu(z)$ itself is still to be done (Chapter~\ref{ch:latetime}). \prediction
\end{enumerate}

\section{Status}
Stands: electroweak symmetry breaking, a smooth crossover at $T_c=159.5$\,GeV and $t\approx9\times10^{-12}$\,s, gives the $W$, $Z$ and charged fermions
their masses and puts them on timelike worldlines; the photon stays massless under the unbroken $U(1)$ and outside the matter sector. \derived{} The
virial partition of IAM's Law applies to the $1/r$ interactions; hadrons are virialised in the Cornell potential and the weak interaction binds
nothing. \derived{} CP violation enters through the baryon asymmetry, whose origin is unknown; the thermodynamic arrow does not need it. Predicted:
$\Sigma=1$ and $\mu_0=-0.136$, tested by Euclid. \prediction{} Open: how much a hadron writes; how the criterion applies to plasma quasiparticles; whether
any physical outcome is selected at the electroweak epoch. \openprob
